\documentclass[a4paper,12pt]{article}
\usepackage[utf8]{inputenc}
\usepackage{listings}
\usepackage{xcolor}
\usepackage{hyperref}

% Farben für Code-Highlighting definieren
\definecolor{codegray}{gray}{0.9}
\definecolor{commentgreen}{rgb}{0,0.6,0}

\lstdefinestyle{pythonstyle}{
    backgroundcolor=\color{codegray},
    commentstyle=\color{commentgreen},
    keywordstyle=\color{blue},
    stringstyle=\color{red},
    basicstyle=\ttfamily\footnotesize,
    breaklines=true,
    frame=single,
    captionpos=b,
    numbers=left,
    numbersep=5pt,
    numberstyle=\tiny\color{gray},
    showspaces=false,
    showstringspaces=false,
    tabsize=4
}


\hypersetup{
    colorlinks=true,
    linkcolor=blue,
    citecolor=blue,
    filecolor=blue,
    urlcolor=blue
}
\parindent 0px
\title{Doodle Jump mit Pygame}
\author{Luca Voithofer}
\date{\today}

\begin{document}

\maketitle

\section{Game Over und Restart-Funktion}
Bisher hatte unser Spiel keine Möglichkeit, nach einem Game Over neu zu starten.  
Dafür fügen wir nun eine Restart-Funktion hinzu, die das Spiel zurücksetzt, wenn die Leertaste (SPACE) gedrückt wird.

\subsection{Game Over Status hinzufügen}
Zunächst benötigen wir eine Variable, die speichert, ob das Spiel vorbei ist.

\begin{lstlisting}[style=pythonstyle, language=Python]
# Game Over Text
gameOverFont = pygame.font.Font("freesansbold.ttf", 64)
gameOverState = False
\end{lstlisting}

\begin{itemize}
\item \textbf{gameOverState:} Speichert, ob das Spiel vorbei ist (True bedeutet Game Over).
\end{itemize} 


\subsection{Game Reset-Funktion}
Damit das Spiel von vorne beginnt, setzen wir alle Variablen auf ihre Anfangswerte zurück.

\begin{lstlisting}[style=pythonstyle, language=Python]
#Reset Game
def resetGame():
    global playerX, playerY, changeX, enemyX, enemyY, changeXenemy, changeYenemy
    global NumberOfEnemies, bulletX, bulletY, changeXbullet, changeYbullet
    global bulletState, gameOverState, score

    # Reset playerposition
    playerX = 370
    playerY = 480
    changeX = 0

    # Replace enemies
    enemyX = []
    enemyY = []
    changeXenemy = []
    changeYenemy = []
    NumberOfEnemies = 8 

    for i in range(NumberOfEnemies):
        enemyX.append(random.randint(0,736))
        enemyY.append(random.randint(50,150))
        changeXenemy.append(0.3)
        changeYenemy.append(40)

    # Reset gamneOverState
    gameOverState = False

    # Reset bullet
    bulletX = 0
    bulletY = 480
    changeXbullet = 0
    changeYbullet = 0.5
    bulletState = "ready"

    # Reset score
    score = 0
\end{lstlisting}

\begin{itemize}
\item \textbf{Alle Spielvariablen} werden auf ihre Ausgangswerte zurückgesetzt.  
\item \textbf{Gegner} werden neu platziert diesmal mit 8 statt 6 Aliens.  
\item \textbf{gameOverState} wird wieder auf False gesetzt, sodass das Spiel neu startet.  
\item \textbf{Score und Kugelzustand} werden zurückgesetzt, damit die Punkte von null beginnen.
\end{itemize}


\subsection{Restart einleiten}
Nun erlauben wir dem Spieler, das Spiel durch Drücken von SPACE neuzustarten, wenn sich das Spiel im GameOverState befindet.

\begin{lstlisting}[style=pythonstyle, language=Python]
# Neustart nach Game Over
if event.key == pygame.K_SPACE and gameOverState == True: 
    resetGame()
\end{lstlisting}

\section{Highscore hinzufügen}
Ein Highscore-System speichert die höchste erreichte Punktzahl und zeigt sie an.  
Dazu fügen wir eine \textbf{Highscore-Variable}, \textbf{eine Anzeige-Funktion} und die \textbf{Aktualisierung der Highscore-Werte} hinzu.

\subsection{Highscore-Variable}
Zunächst definieren wir die Variable highscore, die den höchsten erreichten Punktestand speichert.

\begin{lstlisting}[style=pythonstyle, language=Python]
# Highscore-Variable
highscore = 0
\end{lstlisting}

\begin{itemize}
\item \textbf{highscore = 0:} Initialisiert die Highscore-Variable mit 0.
\end{itemize}


\subsection{Highscore-Anzeige}
Nun erstellen wir eine Funktion, die den Highscore auf dem Bildschirm anzeigt.

\begin{lstlisting}[style=pythonstyle, language=Python]
# Highscore anzeigen
def showHighScore(x, y):
    scorePicture = font.render("HighScore: " + str(highscore), True, (255,255,255))
    screen.blit(scorePicture, (x, y))
\end{lstlisting}

\begin{itemize}
\item \textbf{font.render(...):} Erstellt ein Textbild für den Highscore.  
\item \textbf{screen.blit(...):} Zeichnet den Highscore-Text auf den Bildschirm.  
\item \textbf{x, y:} Definieren die Position der Anzeige.

\end{itemize}  

\subsection{Highscore in der Game Loop aktualisieren}
Die Highscore-Variable muss aktualisiert werden, wenn der aktuelle Score höher ist als der bisherige Highscore.

\begin{lstlisting}[style=pythonstyle, language=Python]
# Highscore aktualisieren
if score > highscore:
    highscore = score
\end{lstlisting}

\begin{itemize}
\item Falls der aktuelle Punktestand (score) größer als der gespeicherte highscore ist, wird der Highscore aktualisiert.
\end{itemize}


\subsection{Highscore auf dem Bildschirm anzeigen}
Nachdem der Highscore aktualisiert wurde, muss die Anzeige-Funktion aufgerufen werden.

\begin{lstlisting}[style=pythonstyle, language=Python]
# Highscore anzeigen
showHighScore(10, 45)
\end{lstlisting}

\begin{itemize}
\item \textbf{showHighScore(10, 45):} es wird die showHighScore-Funktion aufgerufen und der Highscore wird auf die Koordinaten (10, 45) gesetzt
\end{itemize}

\end{document}
